\clearpage\begingroup
\SourceLeaf{157}\HistoricalHead{152}{L}{PHYSIK.\quad BRIEFWECHSEL.}
\begin{center}
[11.]\par\medskip
\textbf{\textsc{Gauss} an \textsc{Gerling}.\quad Göttingen, 19. Januar 1840.}
\end{center}
\noindent . . . . . . . . . . . . . . . . . . . . .\par
\noindent\hspace*{\parindent}In einem Punkte thun Sie meines Erachtens H[errn] \textsc{März} Unrecht, ich meine wegen der Blendgläser. Bei voller Öffnung und hochstehender Sonne springt jedes Glas; auch mir zersprang ein Blendglas am P[assage-]I[nstrument gleich in der ersten Zeit; ich überzeugte mich aber bald, dass dies gar nicht anders sein kann, und machte mir daher gleich selbst Blendkappen von feinem Papier für P. I. und M[ittags-]Kr[eis]. Ich finde die Theorie jenes Vorgangs nirgends berührt, und es scheint, dass manche Physiker davon ganz falsche Vorstellungen haben, da, wenn ich nicht irre \textsc{Munke} (oder war es \textsc{Littrow}?) [es] vor einigen Jahren als etwas Miraculöses betrachtete, dass die Spinnfäden im Brennpunkt eines P. I. nicht verbrennen. Das ist aber ganz natürlich. In der That, ist $m$ die Vergrösserung eines Fernrohrs, die Brennweite $k$mal grösser als der Durchmesser des Objectivs, und setzt man den Sinus des $\odot$ Durchmessers $=\frac1{110}$, so wird, abstrahirt von dem Lichtverlust beim Durchgang durch die Gläser, die Intensität des Sonnenlichtes auf einem Flächenelement sein:
\begin{center}
\begin{tabular}{@{}l@{\;}l@{\quad}p{.76\linewidth}@{}}
1)&$=mm$&vor dem Ocular oder am Platze des Blendglases,\\[4pt]
2)&$=\left(\dfrac{110}{k}\right)^2$&im Brennpunkte oder am Platze der Fäden,
\end{tabular}
\end{center}
\noindent (nemlich die Intensität für einfache Beleuchtung durch die $\odot$ gleich 1 gesetzt).

An Ihrem Fernrohr ist Brennweite 720 Linien, Öffnung, wenn voll, $=43$ Linien, also $k=16\frac34$ und $\frac{110}{k}=6{,}6$. Brauchen Sie also eine 180mahlige Vergrösserung, so ist die erstere Intensität etwa 740mahl grösser als die zweite; für ein gewöhnliches Brennglas ist etwa $k=3$ und auch dann ist die Intensität 1) fast 25mahl stärker. Sie werden, wenn Sie kein Blendglas vorschrauben, und statt dessen ein Stück Zunder vorhalten, diess gleich in Brand gerathen sehen.

Übrigens giebt es gar kein Mittel, jene Intensität zu verändern, wenn man die Vergrösserung beibehält. Aber indem Sie eine Blendkapsel vorstecken, wird der Lichtcylinder vor dem Ocular so viel enger und es wird also von dem Blendglase nur \textls[120]{eine sehr kleine Stelle getroffen}, in welcher aber die Dichtigkeit des Lichtes noch eben so gross ist, wie vorher. . . . . .
\clearpage\endgroup
